\documentclass[11pt,a4paper]{article}

\usepackage[T1]{fontenc}
\usepackage{tgheros}
\renewcommand{\familydefault}{\sfdefault}
\usepackage[margin=2.2cm]{geometry}
\usepackage{microtype}
\usepackage{xcolor}
\usepackage{booktabs}
\usepackage{tabularx}
\usepackage{enumitem}
\usepackage{titlesec}
\usepackage[hidelinks]{hyperref}
\usepackage{parskip}
\usepackage{needspace}

\definecolor{accent}{HTML}{0F766E}
\definecolor{navy}{HTML}{1E3A5F}

\titleformat{\section}{\large\bfseries\color{navy}}{}{0pt}{}
\titleformat{\subsection}{\normalsize\bfseries\color{accent}}{}{0pt}{}
\titlespacing*{\subsection}{0pt}{12pt}{3pt}
\setlist[itemize]{leftmargin=1.2em, itemsep=1.5pt, parsep=0pt, topsep=3pt}

% \project{title}{description}{research question}{data}{approach}{evaluation}{why it's distinct}
\newcommand{\project}[7]{%
  \needspace{14\baselineskip}
  \subsection{#1}
  #2
  \begin{itemize}
    \item \textbf{Research question:} #3
    \item \textbf{Data:} #4
    \item \textbf{Approach:} #5
    \item \textbf{Evaluation:} #6
    \item \textbf{Why it's distinct:} #7
  \end{itemize}}

\pagestyle{plain}

\begin{document}

\begin{center}
  {\LARGE\bfseries\color{navy} Group 10 Project Shortlist}\\[4pt]
  {\small DSL504 / DS504 · Natural Language Processing · IIT Bhilai · Jul--Nov 2026}\\[2pt]
  {\small Agastya Nath · Y.\ Rahul Dev Reddy · Paritosh Lahre · Varun Rao}
\end{center}

\vspace{4pt}
Nine ideas, ranked. None of them overlaps with the topics the other groups chose
(RAG/fact-checking, sentiment, LLM efficiency, NL$\to$SQL, prompt injection).
Each one is scoped for about four weeks of work by four people on free Colab/Kaggle GPUs.

\section*{At a glance}

\renewcommand{\arraystretch}{1.25}
\begin{tabularx}{\textwidth}{@{}c l l >{\raggedright\arraybackslash}X@{}}
\toprule
\textbf{\#} & \textbf{Project} & \textbf{Direction} & \textbf{In one line} \\
\midrule
1 & Do Comments Still Tell the Truth? & Software & Flag code comments made outdated by a code change \\
2 & From Scorecard to Story & Text generation & Accurate cricket commentary from ball-by-ball data \\
3 & Ask, Don't Guess & Software & English$\to$regex that asks clarifying questions \\
4 & When the Law Changes & Legal / India & Predict law sections after the IPC$\to$BNS switch \\
5 & Who's the Doctor? & MT + coreference & Gender bias in English$\to$Hindi translation \\
6 & Parse What You've Never Seen & Software & Log parsing on unseen software systems \\
7 & Did the Minister Answer? & Civic NLP & Evasive answers in Parliament Question Hour \\
8 & Near Hanuman Mandir & Practical / India & Parse messy Indian addresses \\
9 & Unmangling JavaScript & Software & Recover variable names in minified code \\
\bottomrule
\end{tabularx}

\section*{The projects}

\project{1. Do Comments Still Tell the Truth?}
  {When developers change code, they often leave the comment above it untouched, so the comment
   now describes something the code no longer does. For example, a comment says
   ``returns null if not found'', but the code now throws an exception. Given a code change (diff)
   and its existing comment, the system predicts whether the comment has become outdated and
   highlights the stale words.}
  {Does training on artificially generated comment--code mismatches carry over to real ones from
   actual commits?}
  {Panthaplackel et al.\ (AAAI 2021) dataset of real Java inconsistencies; CodeSearchNet
   docstring--code pairs, with the code deliberately changed to create synthetic mismatches.}
  {TF-IDF + logistic regression baseline $\to$ BiGRU with cross-attention $\to$ fine-tuned
   CodeBERT $\to$ few-shot small code LLM.}
  {Precision, recall and F1 on real data; synthetic-only vs.\ real-only training; which kinds of
   artificial change help most.}
  {Group 5 searches for code; this project checks whether existing comments still agree with changed code.}

\project{2. From Scorecard to Story}
  {Generate lively cricket commentary and over/innings summaries from structured ball-by-ball data,
   for example turning ``over 18.3, Bumrah to Buttler, bowled'' into a natural sentence. The main
   problem in generating text from data is fluent text with wrong facts, so a rule-based checker
   recomputes every claim (scores, wickets, milestones) from the data.}
  {Does planning what to say before saying it reduce factual errors, compared with generating
   everything in one step?}
  {Cricsheet (free ball-by-ball data for thousands of matches). It has no commentary text, so
   reference text is built from templates, LLM paraphrases kept only if they pass the checker, and
   about 150 hand-written examples.}
  {Template baseline $\to$ prompted LLM $\to$ fine-tuned T5 $\to$ plan-then-generate T5.}
  {Fact precision and hallucination rate (from the checker), coverage of key events, human
   fluency ratings, diversity.}
  {No group generates text from structured data or uses sports data.}

\project{3. Ask, Don't Guess}
  {Requests like ``match a roll number like 12341550'' are ambiguous. Should any 8 digits match, or
   only ones starting with 1234? Instead of guessing, the system generates several candidate regexes,
   finds a test string they disagree on, and asks the user one question: ``Should 99990000 match?''.
   Regexes can be compared exactly by converting them to automata (DFAs), so this is fully measurable.}
  {How many clarifying questions does it take to reach the intended regex, compared with guessing once?}
  {NL-RX and KB13 (Deep Regex repository), StructuredRegex, plus about 200 descriptions the team
   rewrites to be deliberately vague.}
  {Fine-tuned T5-small with beam search $\to$ group candidates that match exactly the same strings
   $\to$ pick the most informative question $\to$ simulated user that answers automatically.}
  {Exact regex-equivalence accuracy after 0, 1, 2 and 3 questions; whether the correct regex was
   among the candidates at all; a small user study.}
  {Group 4 tests whether NL$\to$SQL handles new combinations of conditions; this project is about
   interaction and ambiguity, and targets regexes instead of SQL.}

\project{4. When the Law Changes}
  {India replaced the Indian Penal Code (IPC) with the Bharatiya Nyaya Sanhita (BNS) on 1 July 2024.
   Sections were renumbered (e.g.\ IPC 420 $\to$ BNS 318), merged, split, or newly added. Every
   existing model that predicts which sections apply to a case's facts was trained on IPC labels.
   This project studies how to move such models to the new code.}
  {Which approach works best: converting predictions afterwards, converting the training labels
   first, or predicting directly from the text of the new BNS sections?}
  {ILSI and IL-TUR (IPC-labelled case facts), the official IPC$\leftrightarrow$BNS mapping table,
   the BNS section texts, and a small hand-labelled BNS test set.}
  {TF-IDF baseline $\to$ InLegalBERT classifier (convert afterwards / convert first) $\to$ zero-shot
   matching of facts to section text $\to$ hybrid.}
  {Micro- and macro-F1, broken down by type of change (renumbered, merged, split, new); top-$k$ recall.}
  {No legal-NLP group. The output is a set of section labels, not a generated answer.}

\project{5. Who's the Doctor?}
  {In Hindi, verbs and adjectives must agree in gender, so translating ``The doctor was tired because
   she worked all night'' forces a choice between feminine and masculine forms. Biased systems pick
   the stereotype (doctor = male) instead of following the pronoun. This project measures how often
   that happens and reduces it.}
  {Can re-ranking beam-search translations using a coreference model reduce gender bias without
   hurting overall translation quality?}
  {WinoMT (about 3,900 English test sentences with known gender), extended with Indian-context
   occupations; FLORES-200 to check general quality.}
  {Translate with IndicTrans2, NLLB, mBART and an LLM; read the Hindi gender automatically via word
   alignment and a morphological tagger; fix the bias with coreference re-ranking and gender tags
   in the source.}
  {Gender accuracy; male--female and stereotype--anti-stereotype accuracy gaps; BLEU/chrF on FLORES.}
  {No group works on machine translation or bias, and it uses coreference, the latest lecture topic.}

\project{6. Parse What You've Never Seen}
  {Software systems write millions of log lines such as
   ``Received block blk\_-160\ldots{} of size 91178 from /10.250.10.6''. Before any monitoring or
   debugging, each line must be split into a fixed template
   (``Received block $\langle*\rangle$ of size $\langle*\rangle$ from $\langle*\rangle$'') and its
   variables. This is a sequence-labelling problem; the challenge is software the model has never seen.}
  {Which methods work on unseen systems, and how many labelled lines does a new system need?}
  {Loghub-2.0 (14 real systems with correct templates) and Loghub-2k (16 systems, 2,000 lines each).}
  {Drain heuristic parser $\to$ BiLSTM-CRF tagger $\to$ fine-tuned RoBERTa tagger $\to$ few-shot
   LLM, evaluated by training on all systems except one and testing on the one left out.}
  {Grouping accuracy, parsing accuracy and their F1 versions; lines per second; a learning curve
   over the number of labelled lines.}
  {No group works on logs or DevOps. This is the lowest-risk option on the list.}

\project{7. Did the Minister Answer?}
  {Lok Sabha questions usually have parts (a), (b), (c), and answers often deflect, e.g.\
   ``(a) to (c): this is a continuous process''. The system splits each question into its parts,
   matches each part to the answer, and labels it answered, partially answered, or not answered.
   It then compares evasion rates across ministries and sessions.}
  {Does transfer from an English political-evasion dataset work on Indian parliamentary answers?}
  {Lok Sabha / Rajya Sabha Q\&A archives; QEvasion (US presidential interviews); about 500
   question parts labelled by the team, with agreement between two annotators checked.}
  {Deflection-phrase rules $\to$ zero-shot natural language inference (NLI) $\to$ DeBERTa transferred
   from QEvasion $\to$ few-shot LLM.}
  {Macro-F1; agreement between annotators (Cohen's $\kappa$); evasion rate by ministry and year.}
  {No civic or political NLP group. It isn't fact-checking: it judges whether the question was
   addressed, not whether the answer is true.}

\project{8. Near Hanuman Mandir}
  {Indian addresses look like ``flat 3b, behind SBI atm nr hanuman mandir, supela bhilai 490023'':
   landmark-based, abbreviated, mixed-language, in no fixed order. The system tags each part (house,
   building, landmark, locality, city, PIN) and standardizes it against the official PIN directory.
   It is useful for delivery, banking KYC and emergency services.}
  {How well do models trained on synthetic messy addresses work on real, hand-written ones?}
  {OpenStreetMap India (public places only), the India Post PIN directory, a noise generator for
   abbreviations, typos and transliteration, and about 300 real test addresses of public places.}
  {Rules + list of place names $\to$ CRF $\to$ BiLSTM-CRF $\to$ fine-tuned MuRIL/XLM-R $\to$
   few-shot LLM; edit-distance matching against the PIN directory.}
  {Entity-level F1 for each field; exact match for the whole address; robustness as noise
   increases; gap between synthetic and real data.}
  {No similar project; it only shares the general technique of sequence labelling.}

\project{9. Unmangling JavaScript}
  {Minified JavaScript replaces meaningful names with single letters:
   \begin{center}\small\texttt{function a(b,c)\{return b.filter(d=>d.price>c)\}}\end{center}
   The system predicts readable names
   (\texttt{items}, \texttt{maxPrice}, \texttt{item}) from how each variable is used. Training
   data is free: minify open-source code and keep the original names as the labels.}
  {Do models understand how a variable is used, or just memorize common names like
   \texttt{i}, \texttt{e}, \texttt{err}?}
  {Self-generated by minifying permissively licensed GitHub code with Terser, split by repository
   so no code leaks between training and test.}
  {Frequency baseline $\to$ n-gram/CRF model $\to$ fine-tuned CodeT5 $\to$ few-shot code LLM.}
  {Exact-name accuracy; subtoken F1 (partial credit, e.g.\ \texttt{userList} vs.\ \texttt{users});
   common vs.\ domain-specific names; human readability ratings.}
  {Group 5 retrieves code from English queries; this project generates names inside the code.}

\section*{Quick picks}
\begin{itemize}
  \item \textbf{Software direction:} \#1 (safest option: \#6)
  \item \textbf{Best presentation:} \#2
  \item \textbf{Most novel:} \#3 (or \#7 if you can start labelling data this week)
  \item \textbf{India-specific relevance:} \#4 or \#5
\end{itemize}

\end{document}
